\section{Kết luận, hạn chế và liên kết}

\subsection{Kết quả đạt được}
\begin{itemize}
\item Phần 1: robot chạy trong Gazebo; \code{/scan}, \code{/odom}, \code{/cmd\_vel} hoạt động (log \file{logs/01\_topic\_check.log}).
\item Phần 2: dựng và lưu bản đồ \file{maps/map.pgm}, \file{maps/map.yaml}; không có tường nhân đôi.
\item Phần 3: nạp bản đồ bằng \code{map\_server}, AMCL hội tụ (sigma khoảng 0,15\,m, $8^\circ$) sau khi đặt initial pose bằng khớp scan--bản đồ.
\item Phần 4: Nav2 đưa robot đến cả 4 goal ở các khu vực khác nhau mà không điều khiển tay; khoảng cách nhỏ nhất tới tường 0,37\,m.
\end{itemize}

\subsection{Hạn chế và những điều cần nói thẳng}
\begin{itemize}
\item \textbf{World tự tạo, nhỏ hơn} world \code{maze} có sẵn; \textbf{cảm biến bị giảm tần số} bằng URDF overlay, và cần \code{LIBGL\_ALWAYS\_SOFTWARE=1}. Đề cho phép dùng môi trường có sẵn nhưng không cấm môi trường tự tạo; lý do là hiệu năng máy.
\item \textbf{Bản đồ có hậu xử lý} bằng \code{clean\_map.py} (xóa điểm lẻ bên ngoài tường); bản thô được giữ lại trong \file{maps/map\_raw.pgm}.
\item \textbf{Bản đồ được quét bằng lái tay}, đúng yêu cầu đề; node tự động bám tường em viết thử không dùng.
\item \textbf{Initial pose lấy từ khớp scan--bản đồ}, nên số liệu ``AMCL trùng với khớp scan'' không phải kiểm chứng độc lập; chưa có vị trí thật từ Gazebo để so (topic ground truth lỗi QoS). Kiểm chứng định tính là tia LiDAR trùng tường và robot tới đúng các goal.
\item \textbf{Goal 1 thất bại lần đầu} (planner không tạo được đường), thành công khi thử lại; chưa xác định chắc nguyên nhân.
\item Số liệu goal ghi bằng script \code{send\_goals.py} (cùng action \code{NavigateToPose}); phần đặt goal bằng công cụ Nav2 Goal của RViz2 được thể hiện trong video/ảnh.
\item Lần gửi goal thứ 5 (để chụp ảnh) thất bại với lỗi TF \emph{Transform data too old} sau khi hệ thống chạy lâu; nguyên nhân gốc chưa rõ (mục~\ref{sec:goal5}).
\item Khi Wi-Fi rớt, Gazebo Transport (multicast) mất kết nối làm \code{/odom}, \code{/scan} ngừng; cần chạy lại mô phỏng. Đây là lỗi môi trường, không phải lỗi thuật toán.
\end{itemize}

\subsection{Liên kết}
\begin{itemize}
\item Video demo (Google Drive, công khai): \texttt{[LINK-VIDEO-GOOGLE-DRIVE]}
\item GitHub repository (Public): \texttt{https://github.com/[GITHUB-USER]/[TÊN-REPO]}
\end{itemize}

\subsection{Tài liệu tham khảo}
\begin{itemize}
\item TurtleBot~4 Simulator: \url{https://github.com/turtlebot/turtlebot4_simulator}
\item TurtleBot~4 User Manual: \url{https://turtlebot.github.io/turtlebot4-user-manual/}
\item SLAM Toolbox: \url{https://github.com/SteveMacenski/slam_toolbox}; Nav2: \url{https://docs.nav2.org}
\end{itemize}
